\chapter{Zmiana miary probabilistycznej i twierdzenie Girsanowa}
\label{ch:girsanov}

\section{Jedna trajektoria, dwa prawa: dlaczego zmieniamy miarę?}

W równaniu $\dd X_t=b(X_t)\,\dd t+\sigma(X_t)\,\dd W_t$
część deterministyczna i losowa odgrywają różne role.
Nasuwa się pytanie, czy można zmienić \emph{prawdopodobieństwa
przypisywane trajektoriom} tak, aby część dryfu została
włączona do nowego procesu Wienera. Odpowiedź daje
\emph{twierdzenie Girsanowa}. Nie zmieniamy tu wartości $X_t(\omega)$
ani samego zbioru trajektorii; zmieniamy miarę, według której
liczymy oczekiwania.

Niech $(\Omega,\calF,(\calF_t)_{0\le t\le T},\Prob)$
będzie filtrowaną przestrzenią probabilistyczną spełniającą
warunki zwykłe. Jeżeli $Z_T>0$ prawie na pewno i
$\E_{\Prob}Z_T=1$, to możemy zdefiniować nową miarę
$\mathbf Q$ na $\calF_T$ przez
$\mathbf Q(A)=\E_{\Prob}[Z_T\ind A]$.
Tak więc dla całkowalnej względem $\mathbf Q$ zmiennej $Y$
\[
\E_{\mathbf Q}Y=\E_{\Prob}(Z_TY).
\]
Gęstość $Z_T$ \emph{przeważa} trajektorie: zdarzenia,
na których $Z_T$ jest duże, otrzymują większą masę.

\begin{lemat}[Warunkowy wzór Bayesa]
\label{lem:bayes-change}
Załóżmy, że $Z_T>0$ $\Prob$-prawie na pewno,
$\E_{\Prob}Z_T=1$, a $\mathbf Q=Z_T\cdot\Prob$
na $\calF_T$. Niech $Z_t=\E_{\Prob}[Z_T\mid\calF_t]$.
Wówczas $Z_t>0$ $\Prob$-prawie na pewno i dla
$0\le t\le T$ oraz $Y\in L^1(\mathbf Q)$
\[
\E_{\mathbf Q}[Y\mid\calF_t]
=\frac{\E_{\Prob}[Z_TY\mid\calF_t]}{Z_t}
\quad\mathbf Q\text{-prawie na pewno}.
\tag{8.1}\label{eq:bayes-measure}
\]
\end{lemat}
\begin{proof}
Ponieważ $Z_T>0$ prawie na pewno, gdyby
$A=\{Z_t=0\}\in\calF_t$ miało dodatnie prawdopodobieństwo,
to $\E[Z_T\ind A]=\E[Z_t\ind A]=0$,
co przeczyłoby dodatniości $Z_T$. Zatem mianownik
jest dodatni prawie na pewno. Oznaczmy prawą stronę
przez $U$. Dla każdego $A\in\calF_t$ mamy
\[
\E_{\mathbf Q}[\ind A U]
=\E_{\Prob}[Z_T\ind A U]
=\E_{\Prob}[Z_t\ind A U]
=\E_{\Prob}[\ind A\E_{\Prob}(Z_TY\mid\calF_t)]
=\E_{\mathbf Q}[\ind A Y].
\]
W drugim przejściu użyliśmy $\calF_t$-mierzalności $U$;
dla nieograniczonych zmiennych najpierw rozpatrujemy
część dodatnią i ujemną albo obcięcia i granicę monotoniczną.
Ponadto $\E_{\mathbf Q}|U|\le\E_{\mathbf Q}|Y|$.
Jest to dokładnie definicja wartości oczekiwanej warunkowej.
\end{proof}

\section{Wykładnik stochastyczny: skąd pochodzi gęstość?}

Niech $W=(W^1,\ldots,W^m)$ będzie $m$-wymiarowym
procesem Wienera, a $\theta=(\theta_t)_{0\le t\le T}$
procesem przewidywalnym o wartościach w $\R^m$.
Na razie załóżmy, że istnieje deterministyczna stała
$K<\infty$ taka, iż $\norm{\theta_t}\le K$
dla $\dd t\otimes\Prob$-prawie wszystkich $(t,\omega)$.
Definiujemy
\[
M_t=-\int_0^t\langle\theta_s,\dd W_s\rangle,
\qquad
Z_t=\exp\left(M_t-\frac12\int_0^t\norm{\theta_s}^2\,\dd s\right).
\tag{8.2}\label{eq:girsanov-density}
\]
Zastosowanie wzoru Itô do $e^{M_t-[M]_t/2}$
pokazuje, że
\[
\dd Z_t=-Z_t\langle\theta_t,\dd W_t\rangle,
\qquad Z_0=1.
\tag{8.3}\label{eq:density-sde}
\]
Poprawka $-\frac12[M]_t$ jest niezbędna: bez niej
wykładnik zwykłej całki martyngałowej ma na ogół
niezerowy dryf.

\begin{lemat}[Wykładnik ograniczonego integrandu jest martyngałem]
\label{lem:bounded-exponential}
Jeśli $\theta$ jest przewidywalny i
$\norm{\theta}\le K$ $\dd t\otimes\Prob$-prawie wszędzie,
to $Z$ z \eqref{eq:girsanov-density} jest ściśle dodatnim,
ciągłym, jednostajnie całkowalnym martyngałem na $[0,T]$.
Dla każdego $p>1$ zachodzi
\[
\sup_{0\le t\le T}\E Z_t^p
\le\exp\left(\frac{p(p-1)}2 K^2T\right),
\tag{8.4}\label{eq:z-moments}
\]
a w szczególności $\E Z_T=1$.
\end{lemat}
\begin{proof}
Najpierw ustalmy $p=2$. Wybierzmy czasy
$\tau_n=\inf\{t\le T:Z_t\ge n\}\wedge T$.
Z wzoru Itô dla $Z^2$ otrzymujemy
\[
Z_{t\wedge\tau_n}^2
=1-2\int_0^{t\wedge\tau_n}Z_s^2\langle\theta_s,\dd W_s\rangle
+\int_0^{t\wedge\tau_n}Z_s^2\norm{\theta_s}^2\,\dd s.
\]
Po zatrzymaniu całka stochastyczna ma zerową średnią,
a $\E Z_{t\wedge\tau_n}^2\le
1+K^2\int_0^t\E Z_{s\wedge\tau_n}^2\,\dd s$.
Gronwall daje oszacowanie jednostajne w $n$.
Ponieważ ciągłe trajektorie $Z$ są ograniczone na
$[0,T]$, mamy $\tau_n=T$ od pewnego losowego $n$;
Fatou i oszacowanie dają $\E Z_t^2\le e^{K^2T}$.
Rodzina zatrzymanych martyngałów $Z^{\tau_n}$
jest zatem jednostajnie całkowalna w każdej ustalonej
chwili (ma jednostajnie ograniczone drugie momenty).
Przechodząc do granicy w warunkowych tożsamościach
martyngałowych otrzymujemy $\E(Z_t\mid\calF_s)=Z_s$.
W szczególności $\E Z_t=1$.

Dla dowolnego $p>1$ zapisujemy
\[
Z_t^p=\underbrace{\exp\left(
-p\int_0^t\langle\theta_s,\dd W_s\rangle
-\frac{p^2}2\int_0^t\norm{\theta_s}^2\,\dd s\right)}_{Z_t^{(p)}}
\exp\left(\frac{p(p-1)}2\int_0^t\norm{\theta_s}^2\,\dd s\right).
\]
Proces $Z^{(p)}$ jest martyngałem na mocy właśnie
udowodnionego przypadku stosowanego do ograniczonego
integrandu $p\theta$; stąd
$\E Z_t^p\le e^{p(p-1)K^2T/2}\E Z_t^{(p)}$.
Wreszcie proces $Z_t=\E[Z_T\mid\calF_t]$ jest
jednostajnie całkowalny jako rodzina wartości oczekiwanych
warunkowych jednej zmiennej całkowalnej.
\end{proof}

\begin{uwaga}[Granica twierdzenia]
Założenie ograniczoności $\theta$ jest wystarczające,
lecz niekonieczne. W dalszej teorii używa się m.in.
warunku Novikova
$\E\exp(\frac12\int_0^T\norm{\theta_s}^2\,\dd s)<\infty$.
W tej iteracji pełny dowód przeprowadzamy przy
ograniczonym integrandzie; nie utożsamiamy lokalnej
martyngałowości wykładnika z zachowaniem jego średniej.
\end{uwaga}

\section{Twierdzenie Girsanowa z pełnym dowodem}

\begin{twierdzenie}[Girsanow przy ograniczonym dryfie]
\label{th:girsanov-bounded}
Niech $T>0$, niech $W$ będzie $m$-wymiarowym
procesem Wienera względem filtracji $(\calF_t)$
spełniającej warunki zwykłe i niech $\theta$ będzie
przewidywalnym procesem o wartościach w $\R^m$,
$\norm{\theta_t}\le K$ dla $\dd t\otimes\Prob$-prawie
wszystkich $(t,\omega)$. Niech $Z$ określa
\eqref{eq:girsanov-density} i niech $\mathbf Q$ na
$\calF_T$ będzie zadana przez $\dd\mathbf Q=Z_T\,\dd\Prob$.
Wówczas proces
\[
\widetilde W_t=W_t+\int_0^t\theta_s\,\dd s,
\qquad 0\le t\le T,
\tag{8.5}\label{eq:shifted-wiener}
\]
jest standardowym $m$-wymiarowym procesem Wienera
względem tej samej filtracji, ale pod miarą $\mathbf Q$.
Ponadto $\Prob$ i $\mathbf Q$ są równoważne na
$\calF_T$ (mają te same zbiory miary zero).
\end{twierdzenie}
\begin{proof}
\emph{Krok 1: zmiana miary jest prawidłowa.}
Z lematu \ref{lem:bounded-exponential} otrzymujemy
$Z_T>0$ i $\E Z_T=1$. Jest więc to gęstość miary
probabilistycznej. Ścisła dodatniość zapewnia
równoważność miar.

\emph{Krok 2: warunkowa funkcja charakterystyczna.}
Ustalmy $u\in\R^m$ i oznaczmy
\[
H_t=\exp\left(i\langle u,\widetilde W_t\rangle
                    +\frac12\norm u^2t\right).
\]
Na mocy wzoru Itô dla rzeczywistej i urojonej
części (albo przez bezpośredni rachunek pochodnych)
$\dd H_t=H_t(i\langle u,\dd W_t\rangle
+i\langle u,\theta_t\rangle\,\dd t)$.
Ze wzoru na iloczyn i \eqref{eq:density-sde} mamy
\[
\dd(Z_tH_t)=Z_tH_t\langle iu-\theta_t,\dd W_t\rangle.
\tag{8.6}\label{eq:girsanov-product}
\]
Wynika to z \emph{dokładnego skasowania} dwóch
członów dryfu: $+iZ_tH_t\langle u,\theta_t\rangle\,\dd t$
i $-iZ_tH_t\langle u,\theta_t\rangle\,\dd t$
pochodzącego z kowariacji.

\emph{Krok 3: lokalny martyngał jest prawdziwy.}
Na $[0,T]$ moduł $H_t$ jest ograniczony przez
$e^{\norm u^2T/2}$. Proces $Z$ ma ograniczone
drugie momenty z \eqref{eq:z-moments}; przy
ograniczonym $\theta$ kwadrat integrandu w
\eqref{eq:girsanov-product} ma skończoną wartość
oczekiwaną po scałkowaniu w czasie. Zatem $ZH$
jest (zespolonym) kwadratowo całkowalnym martyngałem.
Dla $s\le t$ dostajemy
\[
\E_{\Prob}[Z_tH_t\mid\calF_s]=Z_sH_s.
\]
Ponieważ $Z_t=\E_{\Prob}[Z_T\mid\calF_t]$,
wzór Bayesa \eqref{eq:bayes-measure} zastosowany do
$H_t$ daje $\E_{\mathbf Q}[H_t\mid\calF_s]=H_s$.
Po podzieleniu przez dodatni $\calF_s$-mierzalny
czynnik $e^{\norm u^2s/2}e^{i\langle u,\widetilde W_s\rangle}$
otrzymujemy
\[
\E_{\mathbf Q}[e^{i\langle u,\widetilde W_t-\widetilde W_s\rangle}
\mid\calF_s]
=e^{-\norm u^2(t-s)/2}.
\tag{8.7}\label{eq:girsanov-char}
\]
\emph{Krok 4: identyfikacja ruchu Browna.}
Prawa strona jest funkcją charakterystyczną rozkładu
$\Normal(0,(t-s)I_m)$ i nie zależy od $\calF_s$.
Jednoznaczność funkcji charakterystycznych i
argument z klasą monotoniczną dają niezależność
przyrostu od $\calF_s$. Proces $\widetilde W$ ma
ciągłe trajektorie i $\widetilde W_0=0$, zatem
spełnia definicję procesu Wienera pod $\mathbf Q$.
\end{proof}

\begin{uwaga}[Jak odczytywać znak?]
W naszej konwencji gęstość zawiera
$-\int\theta\,\dd W$, a nowy ruch Browna ma postać
$\widetilde W=W+\int\theta\,\dd t$.
Spotyka się także przeciwną konwencję, w której
w wykładniku stoi znak dodatni. Wtedy znak przed
dryfem w nowym procesie Wienera również się zmienia.
Nie ma sprzeczności, ale znaki muszą być uzgodnione.
\end{uwaga}

\section{Jak zmienia się równanie stochastyczne?}

\begin{wniosek}[Transformacja dryfu w SDE]
\label{cor:girsanov-sde}
Załóżmy warunki twierdzenia \ref{th:girsanov-bounded}
i niech proces $X$ spełnia pod $\Prob$, dla
$t\in[0,T]$, równanie
\[
X_t=X_0+\int_0^tb_s\,\dd s+\int_0^t\sigma_s\,\dd W_s,
\]
gdzie $b$ jest progresywnie mierzalny,
$\int_0^T\norm{b_s}\,\dd s<\infty$ prawie na pewno,
a $\sigma$ jest przewidywalny i
$\int_0^T\norm{\sigma_s}_{\mathrm F}^2\,\dd s<\infty$
prawie na pewno. Wtedy pod $\mathbf Q$ ten sam proces $X$
spełnia
\[
X_t=X_0+\int_0^t(b_s-\sigma_s\theta_s)\,\dd s
        +\int_0^t\sigma_s\,\dd\widetilde W_s.
\tag{8.8}\label{eq:girsanov-sde}
\]
\end{wniosek}
\begin{proof}
Tożsamość dla prostych procesów przewidywalnych wynika
z $W_t=\widetilde W_t-\int_0^t\theta_s\,\dd s$
po zastosowaniu liniowości sum całkowych.
Dla ogólnego $\sigma$ wprowadźmy czasy
$\tau_n=\inf\{t:\int_0^t\norm{\sigma_s}_{\mathrm F}^2\,\dd s\ge n\}\wedge T$.
Na przedziale zatrzymanym proces $\sigma$ należy do
$L^2(\dd t\otimes\Prob)$; ponadto zmiana miary może
pozbawić nas globalnego $L^2(\mathbf Q)$, ale przez
zatrzymanie dodatkowo w chwilach, gdy $Z$ lub $1/Z$
opuszczają przedział $[1/n,n]$, otrzymujemy równoważne
normy $L^2$ na zatrzymanym przedziale.
Przybliżając $\sigma$ prostymi procesami w obu normach,
otrzymujemy tożsamość całek najpierw lokalnie, a następnie
na całym $[0,T]$, dzięki wyczerpaniu przedziału przez
czasy lokalizujące. Ponieważ $\theta$ jest ograniczony,
całka $\int_0^t\sigma_s\theta_s\,\dd s$ jest określona
na mocy nierówności Cauchy'ego--Schwarza.
\end{proof}


\begin{uwaga}[Dokładne znaczenie lokalizacji przy zmianie miary]
W dowodzie wniosku \ref{cor:girsanov-sde} trzeba zachować
równoważność norm pod dwiema miarami, ale tylko na
zatrzymanych przedziałach. Dla ustalonego $n$ wybieramy
$\rho_n=\inf\{t\le T:Z_t\notin[1/n,n]\}\wedge T$.
Proces $Z$ jest ciągły i ściśle dodatni; na każdej trajektorii
$\inf_{t\le T}Z_t>0$ i $\sup_{t\le T}Z_t<\infty$,
a zatem $\rho_n=T$ od pewnego losowego $n$.
Jeśli $H$ jest przewidywalny i zeruje się po $\rho_n$,
to na $\calF_{\rho_n}$ zachodzi
$\dd\mathbf Q=Z_{\rho_n}\,\dd\Prob$, skąd
\[
  \frac1n\E_{\Prob}\int_0^{\rho_n}|H_s|^2\,\dd s
 \le \E_{\mathbf Q}\int_0^{\rho_n}|H_s|^2\,\dd s
 \le n\E_{\Prob}\int_0^{\rho_n}|H_s|^2\,\dd s.
\]
Istotne jest zastosowanie gęstości \emph{zatrzymanej}
$Z_{\rho_n}=\E[Z_T\mid\calF_{\rho_n}]$,
a nie gęstości końcowej $Z_T$ bez kontroli.
Przy dodatkowym zatrzymaniu kwadratowego czasu całkowania
powyższe porównanie umożliwia aproksymację tymi samymi
procesami prostymi pod obiema miarami i uzasadnia
podstawienie $\dd W_t=\dd\widetilde W_t-\theta_t\,\dd t$.
\end{uwaga}

\begin{uwaga}[Co pozostaje niezmienione?]
Równoważna zmiana miary nie zmienia własności
ścieżkowych zachodzących prawie na pewno, np. ciągłości
lub wariacji kwadratowej konkretnego procesu.
Zmienia natomiast rozkłady trajektorii i ich wartości
oczekiwane. Wzór \eqref{eq:girsanov-sde} pokazuje,
jak przesunięcie szumu zmienia dryf.
\end{uwaga}

\begin{przyklad}[Wspólny dryf i różne miary: dlaczego horyzont ma znaczenie?]
Dla stałego $c\ne0$ na każdym $[0,T]$ miary prawa
$W$ i $W+ct$ są równoważne, a ich gęstość jest dana
przez wykładnik $e^{cW_T-c^2T/2}$. Nie istnieje jednak
jedna gęstość na całym nieskończonym horyzoncie,
która czyniłaby te prawa równoważnymi na
$C([0,\infty);\R)$: z mocnego prawa wielkich liczb
dla przyrostów w chwilach całkowitych wynika
$W_n/n\to0$ prawie na pewno, podczas gdy
$(W_n+cn)/n\to c$. Zdarzenie
$\{\lim_{n\to\infty}w(n)/n=0\}$ ma więc prawdopodobieństwo
$1$ pod pierwszym prawem, a $0$ pod drugim.
Równoważność na każdym skończonym przedziale nie oznacza
równoważności na całej półprostej.
\end{przyklad}

\section{Zadania rozwiązane i kontrprzykłady intuicyjne}

\begin{zadanie}[Ruch Browna z dryfem jako ruch Browna pod nową miarą]
Niech $W$ będzie procesem Wienera pod $\Prob$, $c\in\R$
i $X_t=W_t+ct$ na $[0,T]$. Skonstruować miarę,
pod którą $X$ jest ruchem Browna, i porównać
prawdopodobieństwo zdarzenia $\{X_T>a\}$.

\emph{Rozwiązanie.} Wybieramy $\theta_t=c$ i
\[
Z_T=\exp(-cW_T-\tfrac12c^2T),\qquad
\dd\mathbf Q=Z_T\,\dd\Prob.
\]
Z twierdzenia Girsanowa $\widetilde W_t=W_t+ct=X_t$
jest procesem Wienera pod $\mathbf Q$. Stąd
\[
\Prob(X_T>a)=1-\Phi\left(\frac{a-cT}{\sqrt T}\right),
\qquad
\mathbf Q(X_T>a)=1-\Phi\left(\frac a{\sqrt T}\right).
\]
Są to różne liczby mimo tego, że $X_t(\omega)$
jest dosłownie tą samą funkcją na obu przestrzeniach
probabilistycznych. Na tym polega zmiana miary, a nie
„usunięcie dryfu z trajektorii”.
\end{zadanie}

\begin{zadanie}[Usuwanie ograniczonego dryfu przez zmianę miary]
Niech $b:\R\to\R$ będzie ograniczoną funkcją Lipschitzowską,
$\sigma_0>0$ i niech $X$ spełnia
$\dd X_t=b(X_t)\,\dd t+\sigma_0\,\dd W_t$, $X_0=x$.
Znaleźć miarę $\mathbf Q$, pod którą $X_t-x$
jest ruchem Browna przeskalowanym przez $\sigma_0$.

\emph{Rozwiązanie.} Chcemy, aby dryf w
\eqref{eq:girsanov-sde} zniknął, więc ustawiamy
$\theta_t=b(X_t)/\sigma_0$.
Proces $\theta$ jest przewidywalny i ograniczony,
ponieważ $X$ jest ciągły i adaptowany. Definiujemy
$Z_T$ przez \eqref{eq:girsanov-density}. Wówczas
$\widetilde W=W+\int\theta\,\dd t$ jest procesem
Wienera pod $\mathbf Q$, a
$\dd X_t=(b(X_t)-\sigma_0\theta_t)\,\dd t
+\sigma_0\,\dd\widetilde W_t=\sigma_0\,\dd\widetilde W_t$.
Stąd $X_t=x+\sigma_0\widetilde W_t$ jako procesy
na całym $[0,T]$. Jeśli dryf $b$ nie byłby ograniczony,
nie moglibyśmy bez dalszego sprawdzenia zastosować
udowodnionej wersji twierdzenia.
\end{zadanie}

\begin{zadanie}[Wartość oczekiwana pod inną miarą]
Niech $W$ będzie jednowymiarowym procesem Wienera,
a $\mathbf Q$ będzie określona przez
$Z_T=e^{-cW_T-c^2T/2}$, gdzie $c\in\R$.
Obliczyć $\E_{\mathbf Q}W_T$ i sprawdzić wynik
dwiema metodami.

\emph{Rozwiązanie.} Ponieważ $\widetilde W_t=W_t+ct$
jest ruchem Browna pod $\mathbf Q$,
$\E_{\mathbf Q}W_T=-cT$.
Z drugiej strony
$\E_{\mathbf Q}W_T=\E_{\Prob}(W_Te^{-cW_T-c^2T/2})$.
Dla $W_T\sim\Normal(0,T)$ różniczkujemy funkcję
tworzącą momenty $\E e^{uW_T}=e^{u^2T/2}$:
$\E(W_Te^{uW_T})=uT e^{u^2T/2}$.
Po podstawieniu $u=-c$ otrzymujemy
$(-cT)e^{c^2T/2}e^{-c^2T/2}=-cT$.
Dwie metody dają ten sam wynik, a druga pokazuje
jak w praktyce wykorzystuje się gęstość Radona--Nikodyma.
\end{zadanie}

\begin{przyklad}[Dlaczego zmiana miary musi zależeć od całej filtracji?]
Jeżeli $\theta_t$ zależy od $W_s$ dla $s>t$, to nie
jest procesem przewidywalnym i całka Itô
$\int\theta\,\dd W$ w ogólności nie ma sensu w
naszej konstrukcji. Próba „odwrócenia znaku dryfu”
na podstawie przyszłego wyniku nie jest poprawną
strategią Girsanowa. Konieczna jest adaptacja
do bieżącej informacji.
\end{przyklad}

\begin{center}
\begin{tikzpicture}[>=Latex,node distance=8mm]
\node[draw,rounded corners,fill=blue!6,minimum width=35mm,align=center] (p)
{$\Prob$\\$W$ jest procesem Wienera};
\node[draw,rounded corners,fill=orange!8,minimum width=36mm,align=center,right=18mm of p] (q)
{$\mathbf Q=Z_T\cdot\Prob$\\$\widetilde W=W+\int\theta\,\dd t$ jest procesem Wienera};
\draw[->,very thick] (p)--node[above,align=center,scale=.85]
{ważenie trajektorii\\przez $Z_T$}(q);
\end{tikzpicture}
\par\smallskip\small Zmiana miary przekształca prawo szumu, nie jego realizacje.
\end{center}

\section{Co dalej?}
Udowodniona wersja twierdzenia Girsanowa ma przejrzysty
zakres: ograniczony przewidywalny dryf, skończony horyzont,
ciągły szum Wienera. Wersje z integrandami nieograniczonymi,
warunkami Novikova i Kazamakiego oraz z miarami losowymi
skoków wymagają odrębnych twierdzeń.
W kolejnych rozdziałach wrócimy do zmiany miary przy
rozważaniu przestrzeni trajektorii i twierdzenia
Camerona--Martina, a później przy procesach Lévy'ego
oraz równaniach ze skokami.

% ===========================================================================
% ROZDZIAL 9 — Miara Wienera i Cameron--Martin
% ===========================================================================
